\section{The inset: regularity, deficit transfer, level selection}
\label{sec:inset}

Throughout this section $\Om\subset\R^2$ is open with $|\Om|=\pi$, $u\in
H^1_0(\Om)$ is its torsion function ($-\Delta u=1$ in $\Om$, $u=0$ on
$\partial\Om$), $m:=\max_\Om u$, and $\dT=\dT(\Om)=\tfrac\pi8-T(\Om)\ge0$ is
the Saint--Venant deficit. We write $E_v:=\{u>v\}$ for the open super-level
set, $\mu(v):=|E_v|$, $\Per(v):=\Per(E_v)$, and recall the two level
deficits of \cref{def:deficits},
\[
   D_H(v)=\mu(v)\bigl(-\mu'(v)\bigr)-\Per(v)^2\ \ge0,
   \qquad
   D_I(v)=\Per(v)^2-4\pi\mu(v)\ \ge0
\]
(\cref{lem:nonneg}), the level-set deficit identity
$\int_0^m(D_H+D_I)\,dv=4\pi\dT$ (\cref{thm:level-identity}), the Talenti
profile $\mu_B(v):=\pi(1-4v)_+$ with $\mu\le\mu_B$ and the profile identity
$\int_0^{1/4}(\mu_B-\mu)\,dv=\dT$ (\cref{thm:talenti}, \cref{lem:profile}),
the bound $-\mu'(v)\ge4\pi$ a.e.\ (\cref{lem:muprime}), and the flux/coarea
identities $\int_{\{u=v\}}|\nabla u|\,d\HH^1=\mu(v)$,
$-\mu'(v)=\int_{\{u=v\}}|\nabla u|^{-1}\,d\HH^1$ for a.e.\ regular $v$
(\cref{thm:coarea}, \cref{lem:dist-deriv}). By Sard's theorem
(\cref{thm:sard}) and \cref{lem:regular-values} the non-regular values of
$u$ form a Lebesgue-null set, and for every regular $v\in(0,m)$ the set
$E_v$ has finite perimeter with $\HH^1$-rectifiable boundary
$\{u=v\}$ (\cref{thm:torsion-reg}). Absolute constants are denoted $C$ (value
may change) or named ($A_0$, $\delta_\star$); $\delta_\star>0$ is a small
absolute constant, fixed at the end of the section. We set $\tau:=\sqrt\dT$
only at the very end; until then $\tau$ is a free parameter in $(0,\tfrac1{16})$.

\subsection{Standing measure-theoretic convention and interior containment}
\label{ss:containment}

The torsion function $u$ is regular only in the \emph{interior} of $\Om$: by
\cref{thm:torsion-reg}, $u\in C^\infty(\Om)$ (indeed real-analytic) and
$u>0$ in $\Om$, but for a general bounded open $\Om$ we do \emph{not} assume
$u\in C(\overline\Om)$, nor that $u=0$ holds pointwise on $\partial\Om$, nor
that the super-level sets $E_v=\{u>v\}$ are compactly contained in $\Om$.
Indeed they need not be: for the punctured disc $\Om=B_1\setminus\set0$
(area $\pi$) the torsion function is $u=u_{B_1}=(1-|x|^2)/4$ (the removed
point $\set0$ is polar, hence invisible to $H^1_0$), and every super-level
set $\set{u\ge t}=\overline{B_{\sqrt{1-4t}}}$ accumulates at the boundary
point $0\in\partial\Om$, so its closure is \emph{not} a compact subset of
$\Om$. The naive statement ``$\overline{E_v}$ is compactly contained in
$\Om$'' is therefore false in general, and we adopt throughout the only
correct, measure-theoretic, stance: all quantities of the paper ---
$|\Om|$, $T(\Om)$, the asymmetry $\AF$, the level sets $E_v$, their
measures $\mu(v)$ and perimeters $\Per(v)$ --- depend on $\Om$ only through
its measure-theoretic structure, and a \emph{polar} (capacity-zero) subset
of $\partial\Om$ may be added to or removed from $\Om$ without changing any
of them. We make this precise and isolate the containment property we
actually use.

\begin{lemma}[Interior containment of super-level sets, modulo a polar set]
\label{lem:containment}
Let $\Om\subset\R^2$ be bounded open with torsion function $u$, and let
$t>0$. Then:
\begin{enumerate}[label=\textup{(\roman*)},leftmargin=2.4em]
\item\label{it:cont-fp} for a.e.\ $t>0$ the set $E_t=\set{u>t}$ has finite
perimeter, $\Per(E_t)=\HH^1(\set{u=t})$, and $E_t$ is an open subset of
$\Om$ that coincides with its own measure-theoretic interior up to an
$\HH^1$-null set;
\item\label{it:cont-far} for every regular $t>0$ and every connected
component $E$ of $E_t$, the relative boundary of $E$ in $\Om$ lies in the
level set,
\[
   \partial E\cap\Om\ \subset\ \set{u=t}\ \subset\ \Om,
\]
an $\HH^1$-rectifiable curve contained in the \emph{open} set $\Om$ and
carrying the full perimeter, $\Per(E)=\HH^1(\partial^*E)=\HH^1(\partial E\cap
\Om)$; the set $E$ is bounded \textup{(}so $\overline E\subset\R^2$ is
compact\textup{)}, and the $\R^2$-closure $\overline E$ may meet
$\partial\Om$, but $\overline E\cap\partial\Om$ carries \emph{no} perimeter
of $E$: $\HH^1(\partial^*E\cap\partial\Om)=0$. \textup{(}This is the precise,
measure-theoretic substitute for ``$\overline E$ compactly contained in
$\Om$'', which fails for the punctured disc.\textup{)}
\item\label{it:cont-polar} consequently, replacing $\Om$ by
$\Om\cup Z$ or $\Om\setminus Z$ for any polar set $Z$ \textup{(}capacity
zero, hence Lebesgue-null and not charged by the perimeter measure of any
finite-perimeter set, \cite[Thm.~4.5, \S4.7]{EvansGariepy2015}\textup{)}
leaves $u$, $T(\Om)$, every $E_t$, $\mu(t)$, $\Per(t)$, $\dT$ and $\AF(\Om)$
unchanged. In particular the symmetric differences $|\Om\triangle E_t|$ and
$|\Om\triangle G|$ used downstream are insensitive to such $Z$.
\end{enumerate}
\end{lemma}

\begin{proof}
\emph{\ref{it:cont-fp}.} By the coarea formula (\cref{thm:coarea}) applied
to $u\in H^1(\Om)\cap C^\infty(\Om)$, $\int_0^m\Per(E_t)\,dt
=\int_\Om|\nabla u|<\infty$, so $E_t$ has finite perimeter for a.e.\ $t$,
with $\Per(E_t)=\HH^1(\set{u=t})$ at every regular $t$
(\cref{lem:regular-values}). The set $E_t$ is open as a super-level set of
the continuous function $u$ on the open set $\Om$; for a regular value
$\nabla u\neq0$ on $\set{u=t}$, so $\set{u=t}$ is an $\HH^1$-rectifiable
$C^1$ curve, of zero Lebesgue measure, whence $E_t$ equals its
measure-theoretic interior up to that null curve.

\emph{\ref{it:cont-far}.} Fix a regular $t$ and a component $E$ of $E_t$.
$E$ is open (super-level set of the continuous $u$) and bounded ($E\subset
\Om$, $\Om$ bounded), so $\overline E\subset\R^2$ is compact. The relative
boundary $\partial E\cap\Om$ consists of points $q\in\Om$ that are limits of
points of $E$ (where $u>t$) and of points of $\Om\setminus E$; by continuity
of $u$ on the open set $\Om$, $u(q)=t$, so $\partial E\cap\Om\subset
\set{u=t}$, which at the regular value $t$ is an $\HH^1$-rectifiable $C^1$
curve in $\Om$ along which $\nabla u\neq0$ (\cref{lem:regular-values}). The
reduced boundary satisfies $\partial^*E\subset\partial E$, and the part of
$\partial^*E$ inside $\Om$ is $\HH^1$-equivalent to $\set{u=t}\cap\partial E$
(\cref{thm:coarea}, $\Per(E)=\HH^1(\set{u=t}\cap\partial E)$). It remains to
see that $\overline E$ may touch $\partial\Om$ only off the reduced boundary.
Indeed $\partial^*E\subset\partial E$, and the part of $\partial E$ on
$\partial\Om$ is disjoint from $\Om\supset\set{u=t}$; since
$\partial^*E\cap\Om\subset\set{u=t}\subset\Om$ already carries the full
perimeter $\Per(E)=\HH^1(\partial^*E)$, the remainder $\partial^*E\cap
\partial\Om$ has $\HH^1$-measure zero. (For the punctured disc, $\overline E
=\overline{B_{\sqrt{1-4t}}}$ touches $\partial\Om$ at the polar puncture
$0$, which is a single point, hence $\HH^1$-null and not a reduced-boundary
point.) This proves \ref{it:cont-far}.

\emph{\ref{it:cont-polar}.} A polar set $Z\subset\R^2$ has zero
$2$-capacity, hence zero Lebesgue measure and is removable for $H^1_0$: the
torsion function of $\Om$ and of $\Om\cup Z$ coincide ($H^1_0(\Om)=
H^1_0(\Om\cup Z)$ when $Z$ is polar of finite measure, since capacity-zero
sets are negligible for $H^1$ traces, \cite[\S4.7]{EvansGariepy2015}). All
listed quantities are integrals against Lebesgue measure of $u$ or
indicators of finite-perimeter sets, or perimeters thereof; Lebesgue measure
and the perimeter measure (which is absolutely continuous with respect to
$\HH^1\restr{\partial^*E}$ and assigns no mass to the $\HH^1$-null, indeed
Lebesgue-null, polar set $Z$, \cite[\S5.7]{EvansGariepy2015},
\cite[Thm.~16.2]{Maggi2012}) both ignore $Z$. The claim follows.
\end{proof}

\begin{remark}[The role of \cref{lem:containment}]\label{rem:containment}
\cref{lem:containment} is the honest substitute for the (false in general)
assertion that $\overline{E_v}$ is compactly contained in $\Om$. What the
later sections genuinely need is only: \textup{(a)} $E_v$ and each of its
components are open subsets of $\Om$ on which $u$ is real-analytic with
$\nabla u\neq0$ along $\set{u=v}$ at regular $v$ (so $\partial E$ is an
$\HH^1$-rectifiable curve in the \emph{open} set $\Om$, with $\Per(E)=
\HH^1(\partial^*E)$); and \textup{(b)} the measure-theoretic insensitivity
to polar sets, \ref{it:cont-polar}, which guarantees that all symmetric
differences and deficits are well posed. In \cref{sec:annulus} we work with
the \emph{measure-theoretic representative} of $\Etil$ \textup{(}equivalently
the finely-open representative\textup{)}: the topological boundary
$\partial\Etil$ there is the boundary of this representative, which differs
from the reduced and essential boundaries only on an $\HH^1$-null set and
lies in the open set $\Om$. No regularity of $\partial\Om$ is used anywhere.
\end{remark}

\subsection{No interior holes}

We now record a topological regularity property of the super-level sets,
which is a direct consequence of the superharmonicity of $u$
($-\Delta u=1\ge0$) and the minimum principle. We use only the
\emph{interior} regularity of \cref{thm:torsion-reg} --- $u\in C^\infty(\Om)$,
$u>0$ in $\Om$ --- together with \cref{lem:containment}; in particular we
do \emph{not} assume $u\in C(\overline\Om)$ or any regularity of
$\partial\Om$. The phrase ``$\overline K\cap\partial\Om\neq\emptyset$''
below means that $K$, a relatively closed bounded subset of $\Om$, is
\emph{not} compactly contained in $\Om$: its closure in $\R^2$ meets
$\partial\Om$ (equivalently, $K$ has a limit point on $\partial\Om$).

\begin{lemma}[No interior holes]\label{lem:noholes}
For every $v\in(0,m)$ the following hold.
\begin{enumerate}[label=\textup{(\roman*)},leftmargin=2.4em]
\item Every connected component $K$ of the open set
$\{x\in\Om:u(x)<v\}$ has $\overline K\cap\partial\Om\neq\emptyset$.
\item For every \emph{regular} value $v$, each connected component $E$ of
$E_v=\{u>v\}$ has no interior holes: every bounded connected component $H$
of $\R^2\setminus\overline E$ satisfies $\overline H\not\subset\Om$; in
particular $\R^2\setminus\overline E$ is connected up to sets touching
$\partial\Om$, and the inset defined below is simply connected modulo
$\partial\Om$.
\end{enumerate}
\end{lemma}

\begin{proof}
The entire argument takes place in the open set $\Om$, where $u$ is smooth
(\cref{thm:torsion-reg}); no value of $u$ on $\partial\Om$ and no continuity
of $u$ up to $\partial\Om$ is invoked. The negations we contradict are
exactly the statements that the relevant components are \emph{compactly
contained in $\Om$} (closure in $\R^2$ a compact subset of $\Om$), the case
\cref{lem:containment}\ref{it:cont-far} shows cannot occur for super-level
components; here we prove the analogous exclusion for the complementary
sub-level islands and for holes.

\emph{(i).} Suppose, for contradiction, that some component $K$ of
$\{u<v\}$ satisfies $\overline K\Subset\Om$, i.e.\ $\overline K$ is a compact
subset of $\Om$ not meeting $\partial\Om$. We claim $u\equiv v$ on
$\partial K$. Let $x_0\in\partial K$. Since $\overline K\subset\Om$ we have
$x_0\in\Om$, so $u$ is continuous at $x_0$ (interior continuity,
\cref{thm:torsion-reg}; no boundary regularity is used). If $u(x_0)<v$ then by continuity
$u<v$ on a ball around $x_0$, which would be a connected subset of
$\{u<v\}$ meeting $K$, hence contained in $K$, placing $x_0$ in the
\emph{interior} of $K$ --- contradicting $x_0\in\partial K$. If $u(x_0)>v$
then by continuity $u>v$ near $x_0$, so a punctured neighbourhood of $x_0$
is disjoint from $\{u<v\}\supset K$, contradicting $x_0\in\partial K$ (as
$K$ accumulates at $x_0$). Hence $u(x_0)=v$, proving $u=v$ on $\partial K$.

Now $u$ is continuous on the compact set $\overline K$ and attains its
minimum there; since $u<v$ at any interior point of $K$ and $u=v$ on
$\partial K$, the minimum is attained at an interior point and equals
$\min_{\overline K}u<v=\min_{\partial K}u$. But $-\Delta u=1\ge0$ in $\Om$
makes $u$ superharmonic, and the (weak) minimum principle for
superharmonic functions in $C^2$
\cite[Theorem~3.1]{GilbargTrudinger2001} gives
$\min_{\overline K}u=\min_{\partial K}u$. This contradiction proves (i).

\emph{(ii).} Let $v$ be regular and $E$ a component of $E_v$, and suppose
$H$ is a bounded component of $\R^2\setminus\overline E$ with
$\overline H\subset\Om$. Then $H$ is an open subset of $\Om$ with
$\overline H\Subset\Om$, and we consider the two open sets
$W^-:=H\cap\{u<v\}$ and $W^+:=H\cap\{u>v\}$, both contained in $H$, hence
both with closures contained in $\overline H\Subset\Om$.

We first note that every connected component $G$ of $W^-$ (resp.\ of $W^+$)
is in fact a connected component of $\{u<v\}$ (resp.\ of $\{u>v\}$), and is
contained in $H$. Indeed, $G$ is connected and contained in the open set
$\{u\ne v\}$, so it lies in a unique component $\hat G$ of $\{u<v\}$ (resp.\
$\{u>v\}$). If $\hat G$ left $H$ it would, being connected and meeting $H$,
have to cross $\partial H$; but $\partial H\subset\overline E$ and on
$\overline E$ one has $u\ge v$, while points of $\partial H$ with $u>v$ lie in
the open set $E$ (a component of $E_v$, hence open) which is disjoint from
$\hat G$ when $\hat G\subset W^-$, and equals a different component of $E_v$
disjoint from $E$ when $\hat G\subset W^+$ (as $\hat G\cap E\subset W^+\cap
E\subset H\cap E=\emptyset$). In either case $\hat G$ cannot meet $\partial H$,
so $\hat G\subset H$ and $\hat G=G$. Thus each component of $W^-$ is a
connected component of $\{u<v\}$
whose closure lies in $\overline H\Subset\Om$ and so does not meet
$\partial\Om$; by (i) this forces $W^-=\emptyset$. Each connected component
$E'$ of $W^+$ is a connected component of $E_v=\{u>v\}$ with
$\overline{E'}\Subset\Om$; on it $u=v$ on $\partial E'$ and $u>v$ inside, so
$u$ attains an interior maximum strictly above its boundary value. But $u$
is superharmonic ($-\Delta u=1\ge0$), so $-u$ is subharmonic and satisfies
the maximum principle \cite[Theorem~3.1]{GilbargTrudinger2001}:
$\max_{\overline{E'}}u=\max_{\partial E'}u=v$, contradicting $u>v$ inside.
Hence $W^+=\emptyset$ as well.

Therefore $u\equiv v$ on the open set $H$, which has positive Lebesgue
measure; but $v$ is a regular value, so $\{u=v\}$ is an $\HH^1$-rectifiable
curve, of zero Lebesgue measure. This contradicts $|H|>0$, so no such $H$
exists.
\end{proof}

\begin{remark}[Measure-theoretic reading; no boundary regularity]
Part (ii) uses both the minimum principle (via (i), excluding holes that are
sub-level islands) and the maximum principle (excluding super-level islands
separated from $E$), and only interior regularity of $u$
(\cref{thm:torsion-reg}); it is consistent with the standing convention of
\cref{ss:containment}. The statement ``$\overline H\not\subset\Om$'' is read
in the sense of \cref{lem:containment}\ref{it:cont-far}: a putative hole $H$
that is compactly contained in $\Om$ is excluded, while a putative hole that
accumulates only on $\partial\Om$ (as the centre of the punctured disc would,
were $\set0$ removed) is not a genuine interior hole and carries no perimeter
of $E$. This does \emph{not} by itself prove that the full exterior
$\R^2\setminus\overline\Etil$ is connected; that stronger topological input is
supplied unconditionally by the hole-filling construction \cref{lem:fill} of
\cref{ss:fill}, which fills any boundary-touching holes at a cost quadratic in
the isoperimetric deficit. The consequences used unconditionally
here are connectedness of $\Etil$, exclusion of compactly contained holes, and
the measure-theoretic fact $|\{u=\hat v\}|=0$ used repeatedly to identify
symmetric differences. A reader content with these may keep (i) and the line
``$u\equiv v$ on $H$ with $|H|>0$ contradicts $|\{u=v\}|=0$'' as the whole of
(ii).
\end{remark}

\subsection{Deficit transfer to super-level sets}

The torsion function of any super-level set $E_v$ is $u-v$, and the
level-set deficit identity, being homogeneous of degree four, transfers the
scale-invariant deficit from $\Om$ to $E_v$ as a tail of the
positive-integrand identity of \cref{thm:level-identity}.

\begin{lemma}[Deficit transfer]\label{lem:transfer}
For every regular $v\in(0,m)$,
\begin{equation}\label{eq:transfer}
   \dstar(E_v)\ \le\ \frac{\pi^2}{\mu(v)^2}\,\dstar(\Om).
\end{equation}
\end{lemma}

\begin{proof}
On the open set $E_v$ the relevant function is
$w:=(u-v)_+\restr{E_v}=u-v$.  This is the torsion function of $E_v$ in the
weak Sobolev sense.  Indeed, by the capacity characterisation of zero traces,
the quasi-continuous representative of $(u-v)_+$ vanishes outside $E_v$ up to
capacity zero, hence $(u-v)_+\in H^1_0(E_v)$
\cite[\S4.7]{EvansGariepy2015}.  For every
$\varphi\in C_c^\infty(E_v)$,
\[
   \int_{E_v}\nabla w\cdot\nabla\varphi
   =\int_{\Om}\nabla u\cdot\nabla\varphi
   =\int_{\Om}\varphi
   =\int_{E_v}\varphi .
\]
Thus $-\Delta w=1$ in $E_v$ with zero Sobolev trace, and uniqueness in
$H^1_0(E_v)$ gives the claim.  Therefore
$T(E_v)=\int_{E_v}w=\int_{E_v}(u-v)$. Its super-level sets are
$\{w>s\}=\{u>v+s\}=E_{v+s}$ for $s\in(0,m-v)$, so the perimeter and
distribution functions of $w$ at height $s$ coincide with those of $u$ at
height $v+s$; in particular the two level deficits of $w$ at height $s$
equal $D_H(v+s)$ and $D_I(v+s)$.

We now apply the derivation of the level-set deficit identity
(\cref{thm:level-identity}, eq.~\eqref{eq:level-identity}) to $w$ on $E_v$;
the derivation there uses only the flux/coarea identities, the absolute
continuity of the distribution function, and the layer-cake formula, none of
which requires the area to equal $\pi$. For a domain $D$ of area $|D|$ with
torsion function having maximum $m_D$ and level deficits $D_H^D,D_I^D$ it
yields
\[
   \int_0^{m_D}\!\bigl(D_H^D+D_I^D\bigr)\,ds
   \ =\ \tfrac12|D|^2-4\pi T(D)
   \ =\ 4\pi\Bigl(\frac{|D|^2}{8\pi}-T(D)\Bigr)
   \ =\ 4\pi\bigl(T(B_D)-T(D)\bigr),
\]
where $B_D$ is the disc of area $|D|$, $T(B_D)=|D|^2/(8\pi)$; the middle
equality is exactly the computation in the proof of
\cref{thm:level-identity} with $\tfrac12|D|^2$ in place of
$\tfrac12\pi^2$. Taking $D=E_v$, $m_{E_v}=m-v$, and using the identification
of the level deficits of $w$ with $D_H(v+\cdot),D_I(v+\cdot)$,
\[
   4\pi\bigl(T(B_{E_v})-T(E_v)\bigr)
   =\int_0^{m-v}\!\bigl(D_H(v+s)+D_I(v+s)\bigr)\,ds
   =\int_v^m\!\bigl(D_H+D_I\bigr)\,dv'
   \ \le\ \int_0^m\!\bigl(D_H+D_I\bigr)\,dv'
   =4\pi\dT,
\]
the inequality by non-negativity of the integrand (\cref{lem:nonneg}).
Hence $T(B_{E_v})-T(E_v)\le\dT$. Dividing by
$T(B_{E_v})=\mu(v)^2/(8\pi)$ and using $\dstar(\Om)=\tfrac8\pi\dT$, i.e.\
$\dT=\tfrac\pi8\dstar(\Om)$,
\[
   \dstar(E_v)=\frac{T(B_{E_v})-T(E_v)}{T(B_{E_v})}
   \ \le\ \frac{\dT}{\mu(v)^2/(8\pi)}
   =\frac{8\pi\dT}{\mu(v)^2}
   =\frac{8\pi}{\mu(v)^2}\cdot\frac\pi8\,\dstar(\Om)
   =\frac{\pi^2}{\mu(v)^2}\,\dstar(\Om).\qedhere
\]
\end{proof}

The content of \eqref{eq:transfer} is that the \emph{scale-invariant}
torsion deficit of every super-level set is comparable to that of $\Om$
(a tail of a positive integral), even at levels near $\partial\Om$ where the
single-level deficits $D_H(v),D_I(v)$ are individually inflated. This is the
mechanism that carries the small torsion deficit onto the inset.

\subsection{Selection of a good level}

\begin{lemma}[Good level]\label{lem:level}
Let $A_0:=16\pi$. For every $\tau\in(0,\tfrac1{16})$ there is a regular
value $\hat v\in(\tau,2\tau)$ of $u$ with
\begin{equation}\label{eq:goodlevel}
   D_H(\hat v)+D_I(\hat v)\ \le\ A_0\,\frac{\dT}{\tau},
   \qquad
   \mu_B(\hat v)-\mu(\hat v)\ \le\ A_0\,\frac{\dT}{\tau}.
\end{equation}
\end{lemma}

\begin{proof}
Set $I:=(\tau,2\tau)$, of length $|I|=\tau$, and note $I\subset(0,\tfrac18)
\subset(0,\tfrac14)$. By \cref{thm:level-identity} and non-negativity,
$\int_I(D_H+D_I)\,dv\le\int_0^m(D_H+D_I)\,dv=4\pi\dT$. By Chebyshev's
inequality,
\[
   \bigl|\{v\in I:\ D_H(v)+D_I(v)>A_0\dT/\tau\}\bigr|
   \ \le\ \frac{\int_I(D_H+D_I)}{A_0\dT/\tau}
   \ \le\ \frac{4\pi\dT}{A_0\dT/\tau}
   =\frac{4\pi\tau}{A_0}.
\]
By \cref{lem:profile}, $\mu_B-\mu\ge0$ and
$\int_I(\mu_B-\mu)\,dv\le\int_0^{1/4}(\mu_B-\mu)\,dv=\dT$, so again by
Chebyshev,
\[
   \bigl|\{v\in I:\ \mu_B(v)-\mu(v)>A_0\dT/\tau\}\bigr|
   \ \le\ \frac{\dT}{A_0\dT/\tau}
   =\frac{\tau}{A_0}.
\]
With $A_0=16\pi$, the union of these two ``bad'' sets has measure at most
$\dfrac{4\pi\tau}{A_0}+\dfrac{\tau}{A_0}=\dfrac{\tau}{4}+\dfrac{\tau}{16\pi}
<\tau=|I|$. Adding the null set of non-regular values
(\cref{thm:sard}, \cref{lem:regular-values}) does not change this. Hence the
set of regular $\hat v\in I$ satisfying both bounds in \eqref{eq:goodlevel}
has positive measure, and in particular is nonempty.
\end{proof}

From now on $\hat v\in(\tau,2\tau)$ denotes a good regular level as in
\cref{lem:level}, and we write $\hat v<\tfrac14$ (since $\hat v<2\tau<
\tfrac18$), so that $\mu_B(\hat v)=\pi(1-4\hat v)$.

\subsection{Closeness of the super-level set}

\begin{lemma}[Closeness]\label{lem:close}
At the level $\hat v$ of \cref{lem:level},
\begin{equation}\label{eq:close}
   \bigl|\Om\triangle\{u>\hat v\}\bigr|
   \ =\ \pi-\mu(\hat v)
   \ \le\ 8\pi\,\tau+A_0\,\frac{\dT}{\tau}.
\end{equation}
\end{lemma}

\begin{proof}
Since $E_{\hat v}=\{u>\hat v\}\subset\Om$, the symmetric difference is
$\Om\setminus E_{\hat v}=\{x\in\Om:u(x)\le\hat v\}$, of measure
$\pi-\mu(\hat v)$ (the level set $\{u=\hat v\}$ is Lebesgue-null). Split
\[
   \pi-\mu(\hat v)
   =\bigl(\pi-\mu_B(\hat v)\bigr)+\bigl(\mu_B(\hat v)-\mu(\hat v)\bigr).
\]
The first term is $\pi-\pi(1-4\hat v)=4\pi\hat v\le4\pi(2\tau)=8\pi\tau$,
using $\hat v<2\tau$; the second is $\le A_0\dT/\tau$ by
\eqref{eq:goodlevel}. Adding gives \eqref{eq:close}.
\end{proof}

\noindent
We now fix $\delta_\star\in(0,\tfrac1{256}]$ small enough that, whenever
$\dT\le\delta_\star$ and $\tau\in[\sqrt\dT,\tfrac1{16})$ with
$\dT/\tau\le\sqrt\dT$ (e.g.\ at the optimal choice $\tau=\sqrt\dT$, where
the right side of \eqref{eq:close} equals $(8\pi+A_0)\sqrt\dT$), one has
\begin{equation}\label{eq:halfmass}
   8\pi\,\tau+A_0\,\frac{\dT}{\tau}\ \le\ \frac\pi2,
   \qquad\text{hence}\qquad
   \mu(\hat v)\ \ge\ \frac\pi2 .
\end{equation}
Concretely, at $\tau=\sqrt\dT$ this requires $(8\pi+A_0)\sqrt{\delta_\star}
\le\pi/2$, i.e.\ $\delta_\star\le\pi^2/(4(8\pi+A_0)^2)$; with $A_0=16\pi$
this is $\delta_\star\le 1/(2304)$, comfortably compatible with
$\delta_\star\le\tfrac1{256}$ after taking the smaller value. We keep
\eqref{eq:halfmass} as a standing hypothesis for the remainder of the
section (it holds at the optimised scale $\tau=\sqrt\dT$); the value of
$\hat v,\tau$ is otherwise left symbolic and optimised in
\cref{thm:reduction}.

\subsection{One connected inset: the spillover is higher order}

The super-level set $\{u>\hat v\}$ may have several components, but the
total mass of all but the largest is quadratically small in $D_I(\hat v)$,
by a multi-component refinement of the isoperimetric inequality.

\begin{lemma}[Spillover]\label{lem:spill}
Assume \eqref{eq:halfmass}. Let the connected components of $\{u>\hat v\}$
have areas $A_1\ge A_2\ge\cdots>0$, let $\Etil$ be the component of area
$A_1$, and set $S:=\sum_{i\ge2}A_i=\mu(\hat v)-A_1$ for the total spillover.
Then $A_1\ge\pi/4$ and
\begin{equation}\label{eq:spill}
   \bigl|\{u>\hat v\}\setminus\Etil\bigr|
   \ =\ S\ \le\ \frac{D_I(\hat v)^2}{16\pi^3}
   \ \le\ \frac{A_0^2}{16\pi^3}\,\frac{\dT^2}{\tau^2}.
\end{equation}
\end{lemma}

\begin{proof}
\emph{Multi-component isoperimetric inequality.} Perimeter is additive over
the (essentially disjoint) components and the planar isoperimetric
inequality (\cref{thm:isoperimetric}) applies to each, so
$\Per(\hat v)=\sum_i\Per_i\ge\sum_i 2\sqrt{\pi A_i}$ with $\Per_i$ the
perimeter of the $i$-th component. Squaring and using
$\bigl(\sum_i\sqrt{A_i}\bigr)^2=\sum_iA_i+2\sum_{i<j}\sqrt{A_iA_j}$,
\begin{equation}\label{eq:multiiso}
   D_I(\hat v)=\Per(\hat v)^2-4\pi\!\sum_iA_i
   \ \ge\ 4\pi\Bigl[\Bigl(\textstyle\sum_i\sqrt{A_i}\Bigr)^2-\sum_iA_i\Bigr]
   \ =\ 8\pi\!\!\sum_{i<j}\!\sqrt{A_iA_j}.
\end{equation}

\emph{Lower bound on $A_1$.} Since $A_j\le A_1$ we have
$\sqrt{A_j}=A_j/\sqrt{A_j}\ge A_j/\sqrt{A_1}$ for $j\ge2$, hence
$\sum_{j\ge2}\sqrt{A_j}\ge S/\sqrt{A_1}$ and
\[
   \sum_i\sqrt{A_i}\ \ge\ \sqrt{A_1}+\frac{S}{\sqrt{A_1}}
   =\frac{A_1+S}{\sqrt{A_1}}=\frac{\mu(\hat v)}{\sqrt{A_1}}.
\]
Combining with \eqref{eq:multiiso} (in the form $D_I\ge4\pi[(\sum\sqrt{A_i})^2
-\mu(\hat v)]$),
\begin{equation}\label{eq:Aone}
   D_I(\hat v)\ \ge\ 4\pi\Bigl(\frac{\mu(\hat v)^2}{A_1}-\mu(\hat v)\Bigr)
   =4\pi\,\mu(\hat v)\,\frac{\mu(\hat v)-A_1}{A_1}
   =4\pi\,\mu(\hat v)\,\frac{S}{A_1}.
\end{equation}
Suppose, for contradiction, that $A_1<\pi/4$. By \eqref{eq:halfmass},
$\mu(\hat v)\ge\pi/2>A_1$, so there are at least two components and
$S=\mu(\hat v)-A_1>\pi/2-\pi/4=\pi/4$. Then \eqref{eq:Aone} gives
\[
   D_I(\hat v)\ >\ 4\pi\cdot\frac\pi2\cdot\frac{S}{\pi/4}=8\pi\,S>8\pi\cdot\frac\pi4=2\pi^2.
\]
But $D_I(\hat v)\le A_0\dT/\tau\le8\pi\tau+A_0\dT/\tau\le\pi/2<2\pi^2$ by
\eqref{eq:goodlevel} and \eqref{eq:halfmass}, a contradiction. Hence
$A_1\ge\pi/4$.

\emph{Spillover bound.} From \eqref{eq:multiiso},
\[
   D_I(\hat v)\ \ge\ 8\pi\!\sum_{i<j}\!\sqrt{A_iA_j}
   \ \ge\ 8\pi\sqrt{A_1}\sum_{j\ge2}\sqrt{A_j}
   \ \ge\ 8\pi\sqrt{A_1}\,\sqrt{\textstyle\sum_{j\ge2}A_j}
   \ =\ 8\pi\sqrt{A_1\,S},
\]
where the middle step keeps only the pairs $(1,j)$ and the last step uses
$\sum_{j\ge2}\sqrt{A_j}\ge\sqrt{\sum_{j\ge2}A_j}$ (superadditivity of
$t\mapsto\sqrt t$ on non-negative terms). Squaring and using $A_1\ge\pi/4$,
\[
   S\ \le\ \frac{D_I(\hat v)^2}{64\pi^2A_1}
   \ \le\ \frac{D_I(\hat v)^2}{64\pi^2\cdot(\pi/4)}
   =\frac{D_I(\hat v)^2}{16\pi^3}.
\]
The final bound in \eqref{eq:spill} follows from
$D_I(\hat v)\le A_0\dT/\tau$ (\cref{lem:level}).
\end{proof}

\begin{definition}[The inset]\label{def:inset}
Assume $\dT\le\delta_\star$, $\tau$ subject to \eqref{eq:halfmass}, and let
$\hat v\in(\tau,2\tau)$ be a good regular level (\cref{lem:level}). The
\emph{inset} is the connected component $\Etil$ of $\{u>\hat v\}$ of largest
measure, i.e.\ the component of area $A_1$ in \cref{lem:spill}. We write
\[
   \mu_E:=|\Etil|=A_1=\mu(\hat v)-S\ \ge\ \frac\pi4,
   \qquad
   \hat r:=\sqrt{\mu_E/\pi},
   \qquad
   \diso:=\diso(\Etil)=\frac{\Per(\Etil)}{2\sqrt{\pi\mu_E}}-1,
\]
and $\Exc:=\Per(\Etil)-2\pi\hat r=2\pi\hat r\,\diso$.
\end{definition}

By \cref{lem:noholes}\,(ii), $\Etil$ is connected and has no interior hole
\emph{compactly contained in $\Om$}. The stronger \emph{global}
exterior-connectedness ($\R^2\setminus\overline\Etil$ connected) required by
\cref{sec:annulus,sec:fold} is supplied unconditionally by passing to the
hole-filled inset of \cref{lem:fill} below (\cref{ss:fill}); from
\cref{sec:annulus} onward the geometric reduction is run on that saturation,
which inherits all the controls of \cref{prop:insetcontrols} at the same scale.

\subsection{Controls on the inset}

We collect the three controls of the inset used in the subsequent sections.

\begin{proposition}[Inset controls]\label{prop:insetcontrols}
Under \cref{def:inset} there is an absolute constant $C$ such that
\begin{align}
   \diso(\Etil)\ &\le\ \frac{D_I(\hat v)+4\pi S}{8\pi\,\mu_E}
   \ \le\ C\,\frac{\dT}{\tau},
   \label{eq:diso-inset}\\[1mm]
   \dstar(\Etil)\ &\le\ \frac{\pi^2}{\mu_E^2}\,\dstar(\Om)
   +\frac{S^2}{\mu_E^2}
   \ \le\ C\,\dstar(\Om),
   \label{eq:dstar-inset}\\[1mm]
   \bigl|\Om\triangle\Etil\bigr|\ &\le\ 8\pi\,\tau+A_0\frac{\dT}{\tau}
   +\frac{A_0^2}{16\pi^3}\frac{\dT^2}{\tau^2}.
   \label{eq:close-inset}
\end{align}
In particular, with the optimal choice $\tau=\sqrt\dT$ of
\cref{thm:reduction}, $\diso(\Etil)\le C\sqrt\dT$,
$\dstar(\Etil)\le C\dstar(\Om)$ and
$|\Om\triangle\Etil|\le C\sqrt\dT$.
\end{proposition}

\begin{proof}
\emph{\eqref{eq:diso-inset}.} Since $\partial\Etil\subset\{u=\hat v\}$ and
perimeter is additive over components, $\Per(\Etil)\le\Per(\hat v)$, hence
\[
   D_I(\Etil):=\Per(\Etil)^2-4\pi\mu_E
   \ \le\ \Per(\hat v)^2-4\pi\mu_E
   =\bigl(4\pi\mu(\hat v)+D_I(\hat v)\bigr)-4\pi\mu_E
   =D_I(\hat v)+4\pi S,
\]
using $\Per(\hat v)^2=4\pi\mu(\hat v)+D_I(\hat v)$ and
$\mu(\hat v)-\mu_E=S$. From the definition of $\diso$,
$\bigl(1+\diso(\Etil)\bigr)^2=\Per(\Etil)^2/(4\pi\mu_E)=1+D_I(\Etil)/(4\pi\mu_E)$,
so by $\sqrt{1+t}-1\le t/2$,
\[
   \diso(\Etil)\ \le\ \frac{D_I(\Etil)}{8\pi\mu_E}
   \ \le\ \frac{D_I(\hat v)+4\pi S}{8\pi\mu_E}.
\]
By \eqref{eq:spill}, $4\pi S\le 4\pi D_I(\hat v)^2/(16\pi^3)
=D_I(\hat v)^2/(4\pi^2)$; and $D_I(\hat v)\le A_0\dT/\tau\le\pi/2$ by
\eqref{eq:goodlevel}--\eqref{eq:halfmass}, so $4\pi S\le D_I(\hat v)\cdot
\frac{D_I(\hat v)}{4\pi^2}\le\frac{1}{8\pi}D_I(\hat v)\le D_I(\hat v)$.
Hence $D_I(\hat v)+4\pi S\le2D_I(\hat v)\le2A_0\dT/\tau$, and with
$\mu_E\ge\pi/4$,
\[
   \diso(\Etil)\ \le\ \frac{2A_0\dT/\tau}{8\pi\cdot(\pi/4)}
   =\frac{A_0}{\pi^2}\,\frac{\dT}{\tau}\ =:\ C\,\frac{\dT}{\tau}.
\]

\emph{\eqref{eq:dstar-inset}.} We bound $\dstar(\Etil)$ directly, using only
the deficit transfer of \cref{lem:transfer} and the exact additivity of
torsion over connected components; no appeal to \cref{lem:cutdeficit} is
needed (see \cref{rem:fwd}). Let $C_2,C_3,\dots$ be the spillover components
of $E_{\hat v}$, of areas $A_2,A_3,\dots$ with $\sum_{i\ge2}A_i=S$. On each
component $C$ of $E_{\hat v}=\{u>\hat v\}$, the restriction of
$(u-\hat v)_+$ to $C$ belongs to $H^1_0(C)$.  More explicitly, the Sobolev
space on an open set decomposes over its connected components:
if $E=\bigcup_i C_i$ is the disjoint union of open components and
$w\in H^1_0(E)$, then $w_i:=w\mathbf1_{C_i}$ belongs to $H^1_0(C_i)$ and
$\nabla w_i=\mathbf1_{C_i}\nabla w$ a.e.; this follows from the
quasi-continuous trace characterisation of $H^1_0$ and the locality of the
Sobolev gradient.  Applying this to $w=(u-\hat v)_+\in H^1_0(E_{\hat v})$
gives the claim.  Testing against $C_c^\infty(C)$ gives
$-\Delta(u-\hat v)=1$ weakly in $C$, and uniqueness makes this restriction
the torsion function of $C$; hence $T(C)=\int_C(u-\hat v)$. As
$u-\hat v$ is shared and the components are pairwise disjoint, torsion is
\emph{exactly additive}:
\begin{equation}\label{eq:Tadd}
   T(E_{\hat v})=\int_{E_{\hat v}}(u-\hat v)
   =\int_{\Etil}(u-\hat v)+\sum_{i\ge2}\int_{C_i}(u-\hat v)
   =T(\Etil)+\sum_{i\ge2}T(C_i).
\end{equation}
By the Saint--Venant inequality applied to each spillover component
(\cref{prop:scale}, $T(B_E)-T(E)\ge0$, with $T(B_{C_i})=A_i^2/(8\pi)$),
\begin{equation}\label{eq:Tspill}
   \sum_{i\ge2}T(C_i)\ \le\ \sum_{i\ge2}\frac{A_i^2}{8\pi}
   \ \le\ \frac{1}{8\pi}\Bigl(\sum_{i\ge2}A_i\Bigr)^2=\frac{S^2}{8\pi},
\end{equation}
the middle inequality because $\sum A_i^2\le(\sum A_i)^2$ for non-negative
terms. Combining \eqref{eq:Tadd}--\eqref{eq:Tspill} with $\mu_E\le\mu(\hat
v)$,
\[
   T(B_{\Etil})-T(\Etil)
   =\frac{\mu_E^2}{8\pi}-T(E_{\hat v})+\sum_{i\ge2}T(C_i)
   \ \le\ \frac{\mu_E^2-\mu(\hat v)^2}{8\pi}
          +\bigl(T(B_{E_{\hat v}})-T(E_{\hat v})\bigr)+\frac{S^2}{8\pi}
   \ \le\ \dT+\frac{S^2}{8\pi},
\]
where we used $T(B_{\Etil})=\mu_E^2/(8\pi)$, $T(B_{E_{\hat v}})=
\mu(\hat v)^2/(8\pi)$, $\mu_E^2-\mu(\hat v)^2\le0$, and the bound
$T(B_{E_{\hat v}})-T(E_{\hat v})\le\dT$ established in the proof of
\cref{lem:transfer}. Dividing by $T(B_{\Etil})=\mu_E^2/(8\pi)$ and using
$\dT=\tfrac\pi8\dstar(\Om)$,
\begin{equation}\label{eq:dstarE-exact}
   \dstar(\Etil)
   \ \le\ \frac{8\pi\dT}{\mu_E^2}+\frac{S^2}{\mu_E^2}
   \ =\ \frac{\pi^2}{\mu_E^2}\,\dstar(\Om)+\frac{S^2}{\mu_E^2}.
\end{equation}
This is the first inequality in \eqref{eq:dstar-inset} up to the higher-order
term $S^2/\mu_E^2$. The latter is negligible: by \eqref{eq:spill},
$S\le A_0^2\dT^2/(16\pi^3\tau^2)$, and with the standing hypothesis
\eqref{eq:halfmass} (in particular $A_0\dT/\tau\le\pi/2$, i.e.\
$\dT/\tau\le1/32$) one has $S\le\tfrac{A_0^2}{16\pi^3}(\dT/\tau)\cdot\dT
\le\tfrac{A_0^2}{16\pi^3\cdot32}\dT\le\dT$; hence, with $\mu_E\ge\pi/4$,
\[
   \frac{S^2}{\mu_E^2}\ \le\ \frac{S\cdot\dT}{(\pi/4)^2}
   \ \le\ \frac{16}{\pi^2}\,\dT^2
   \ \le\ \frac{16}{\pi^2}\,\dT\cdot\frac\pi8\,\dstar(\Om)
   \ =\ \frac{2}{\pi}\,\dT\,\dstar(\Om)\ \le\ \dstar(\Om),
\]
the last step for $\dT\le\delta_\star\le\pi/2$. Thus, with $\mu_E\ge\pi/4$,
\eqref{eq:dstarE-exact} gives
$\dstar(\Etil)\le\frac{\pi^2}{\mu_E^2}\dstar(\Om)+\dstar(\Om)
\le(16+1)\dstar(\Om)\le C\,\dstar(\Om)$, which is \eqref{eq:dstar-inset}.

\emph{\eqref{eq:close-inset}.} By \cref{lem:close} and \cref{lem:spill},
\[
   |\Om\triangle\Etil|=\bigl(\pi-\mu(\hat v)\bigr)+S
   \ \le\ \Bigl(8\pi\tau+A_0\frac\dT\tau\Bigr)+\frac{A_0^2}{16\pi^3}\frac{\dT^2}{\tau^2}.
\]
For the final assertions, set $\tau=\sqrt\dT$: then $\dT/\tau=\sqrt\dT$, so
$\diso(\Etil)\le C\sqrt\dT$ by \eqref{eq:diso-inset}; the spillover term
$\dT^2/\tau^2=\dT$ is higher order, so $|\Om\triangle\Etil|\le(8\pi+A_0)
\sqrt\dT+C\dT\le C\sqrt\dT$; and \eqref{eq:dstar-inset} is independent of
$\tau$.
\end{proof}

\subsection{Hole-filling: global exterior connectedness at no cost}
\label{ss:fill}

\Cref{lem:noholes} excludes complementary components of $\Etil$ that are
\emph{compactly contained} in $\Om$, but it does not by itself rule out a
bounded complementary component whose closure reaches $\partial\Om$: such a
component evades the interior minimum principle, which is available only away
from $\partial\Om$. The downstream geometry of \cref{sec:annulus,sec:fold},
however, needs the \emph{global} topological fact that the unbounded
complement of the set is its only complement --- the exterior-connectedness
convention of \cref{lem:ext-conn}. We supply this unconditionally by passing
to the \emph{saturation} (hole-filling) of $\Etil$, the standard device of
filling in all bounded complementary components. The decisive quantitative
point is that, by an isoperimetric argument identical in spirit to the
spillover bound \cref{lem:spill}, the filled area is bounded by the
\emph{square} of the isoperimetric deficit, hence is of the same higher order
as the spillover $S$; all inset scales are therefore preserved.

\begin{definition}[Saturation of the inset]\label{def:fill}
The \emph{hole-filled inset} is
\begin{equation}\label{eq:fill-def}
   E^h\ :=\ \Etil\ \cup\ \bigcup\set{K:\ K\text{ a bounded connected
   component of }\R^2\setminus\overline\Etil},
\end{equation}
the union of $\Etil$ with all bounded complementary components of its
$\R^2$-closure. We write $H:=|E^h\setminus\Etil|\ge0$ for the filled area.
\end{definition}

\begin{lemma}[Hole-filling]\label{lem:fill}
Under \cref{def:inset} the hole-filled inset $E^h$ of \cref{def:fill}
satisfies the following.
\begin{enumerate}[label=\textup{(\alph*)},leftmargin=2.4em]
\item\label{it:fill-top} $E^h$ is bounded, open and connected; its complement
$\R^2\setminus\overline{E^h}$ is connected and unbounded \textup{(}it is
exactly the unbounded component of $\R^2\setminus\overline\Etil$\textup{)},
and $\partial E^h\subset\partial\Etil$. Thus $E^h$ is globally hole-free and
satisfies the exterior-connectedness convention \cref{lem:ext-conn}
\textup{(}$\R^2\setminus\overline{E^h}$ connected, with
$\partial(\R^2\setminus\overline{E^h})=\partial E^h$\textup{)}
unconditionally, with no hypothesis on $\partial\Om$.
\item\label{it:fill-per} $\Etil\subset E^h$, $|E^h|=\mu_E+H$, and the
perimeter does not increase:
\begin{equation}\label{eq:fill-per}
   \Per(E^h)\ \le\ \Per(\Etil).
\end{equation}
\item\label{it:fill-key} \textup{(}Key bound.\textup{)} With
$D_I(\Etil)=\Per(\Etil)^2-4\pi\mu_E\ge0$,
\begin{equation}\label{eq:fill-key}
   H\ \le\ \frac{D_I(\Etil)^2}{64\pi^2\,\mu_E}\ \le\ C\,D_I(\Etil)^2 .
\end{equation}
\item\label{it:fill-scales} At the inset scales of
\cref{prop:insetcontrols} \textup{(}$D_I(\Etil)\le C\,\dT/\tau$,
$\mu_E\ge\pi/4$\textup{)} one has $H\le C(\dT/\tau)^2$, so at $\tau=\sqrt\dT$,
$H\le C\,\dT$ \textup{(}higher order, the order of the spillover $S$\textup{)}.
Consequently all the inset controls transfer to $E^h$, with the same rates:
writing $\hat r^h:=\sqrt{|E^h|/\pi}$,
\begin{align}
   |\Om\triangle E^h|\ &\le\ |\Om\triangle\Etil|+H\ \le\ C\,\sqrt\dT,
   &\hat r^h\ &=\ \hat r+O(H)=\hat r+O(\dT);\label{eq:fill-close}\\[1mm]
   \diso(E^h)\ &\le\ \diso(\Etil)\ \le\ C\,\tfrac\dT\tau,
   &\Exc(E^h)\ &\le\ \Exc(\Etil)+C\,H\ \le\ C\,\tfrac\dT\tau;\label{eq:fill-diso}\\[1mm]
   a(E^h)\ &\le\ a(\Etil)+C\,H\ \le\ a(\Etil)+C\,\dT,
   &&\bigl(\text{so }a(E^h)\le C(\tfrac\dT\tau)^{1/2}\bigr);\label{eq:fill-a}\\[1mm]
   \dstar(E^h)\ &\le\ \dstar(\Etil)+\frac{C\,H}{\mu_E}\ \le\ \dstar(\Etil)+C\,\dT
   \ \le\ C\,\dstar(\Om), &&\label{eq:fill-dstar}
\end{align}
where $a(F):=|F\triangle B_{\hat r_F}(z_F)|$ is the ball distance of $F$ at
its optimal centre and $\hat r_F:=\sqrt{|F|/\pi}$.
\end{enumerate}
\end{lemma}

\begin{proof}
\emph{\ref{it:fill-top}.} The open set $\R^2\setminus\overline\Etil$ has, as
the complement of a bounded set, exactly one unbounded connected component
$U_\infty$ (nonempty since $\Etil$ is bounded) and at most countably many
bounded components $\set{K_i}_{i\in I}$, all open and disjoint; by
\eqref{eq:fill-def}, $E^h=\Etil\cup\bigcup_iK_i$, so that as sets
\[
   E^h\ =\ \R^2\setminus\bigl(\partial\Etil\cup U_\infty\bigr),
   \qquad\text{equivalently}\qquad
   \overline{E^h}=\R^2\setminus U_\infty .
\]
We first record the elementary topological facts. Each component $K$ of
$\R^2\setminus\overline\Etil$ (bounded or not) is open with $\partial
K\subset\partial\Etil$: a point $x\in\partial K$ lies neither in the open set
$K$ nor in the open set $\R^2\setminus\overline\Etil\setminus K$ (a union of
other components), so $x\in\overline\Etil$; and $x\notin\Etil$, since $\Etil$
open with $x\in\overline K$ would put a neighbourhood of $x$ inside $\Etil$,
disjoint from $K$. Hence $x\in\partial\Etil$. In particular $\partial
U_\infty\subset\partial\Etil$ and $\partial K_i\subset\partial\Etil$.

\emph{$E^h$ is open.} Since $\Etil$ and each $K_i$ are open, it suffices to
show every $x\in\partial\Etil\cap E^h$ is interior to $E^h$; but in fact no
point of $\partial\Etil$ lies in $E^h$, because $E^h=\Etil\cup\bigcup_iK_i$
is a union of open sets disjoint from $\partial\Etil$
(as $\partial\Etil\cap\Etil=\emptyset$ and $\partial\Etil\cap K_i=\emptyset$,
the latter since $K_i$ is open and $\subset\R^2\setminus\overline\Etil$).
Thus $E^h$ is a union of open sets, hence open. It is bounded since $\Etil$
and all $K_i$ lie in any ball $B$ with $\R^2\setminus B\subset U_\infty$
(such $B$ exists as $U_\infty$ contains the complement of a large ball).

\emph{$E^h$ is connected.} $\Etil$ is connected (\cref{def:inset},
\cref{lem:noholes}). For each $i$, $\partial K_i\subset\partial\Etil$ is
nonempty (a bounded open $K_i\subsetneq\R^2$ has nonempty boundary), so
pick $x\in\partial K_i\subset\partial\Etil$; then $x\in\overline{K_i}\cap
\overline\Etil$, so $\Etil\cup K_i$ is connected (two connected sets with
intersecting closures, both inside the connected $\overline{\Etil\cup K_i}$;
concretely $\overline{K_i}\cap\overline\Etil\ni x$ joins them). All the
$\Etil\cup K_i$ share the connected $\Etil$, so their union $E^h$ is
connected.

\emph{Exterior connectedness.} By the displayed identity
$\overline{E^h}=\R^2\setminus U_\infty$, so $\R^2\setminus\overline{E^h}=
U_\infty$, which is connected and unbounded by construction. Moreover
$\partial E^h=\overline{E^h}\setminus E^h=(\R^2\setminus U_\infty)\setminus
E^h\subset\partial\Etil$ (the points of $\R^2\setminus U_\infty$ not in
$E^h=\Etil\cup\bigcup_iK_i$ lie in $\partial\Etil$, as
$\R^2=\Etil\sqcup\partial\Etil\sqcup(\R^2\setminus\overline\Etil)$ and
$\R^2\setminus\overline\Etil=U_\infty\sqcup\bigcup_iK_i$). Conversely
$\partial U_\infty=\overline{U_\infty}\setminus U_\infty\subset\partial\Etil$
lies in $\partial E^h$ since these points are limits of $U_\infty=\R^2
\setminus\overline{E^h}$ and of $E^h$ (being in $\partial\Etil$, hence limits
of $\Etil\subset E^h$). Thus $\partial(\R^2\setminus\overline{E^h})=\partial
U_\infty=\partial E^h$, and the trichotomy $\R^2=E^h\sqcup\partial E^h\sqcup
(\R^2\setminus\overline{E^h})$ holds. This is exactly the
exterior-connectedness convention of \cref{lem:ext-conn}, established with no
appeal to the minimum principle, hence with no hypothesis on $\partial\Om$.

\emph{\ref{it:fill-per}.} Clearly $\Etil\subset E^h$ and the $K_i$ are
pairwise disjoint and disjoint from $\Etil$, so $|E^h|=\mu_E+\sum_i|K_i|=
\mu_E+H$ with $H=\sum_i|K_i|\ge0$; each $K_i$ is bounded with finite
perimeter $\Per(K_i)<\infty$, since $\Per(\Etil)<\infty$ and
$\partial^*K_i\subset\partial^*\Etil$ $\HH^1$-a.e. (see below).

For the perimeter, we identify the essential boundaries. The plane is
partitioned $\HH^2$-a.e.\ into the measure-theoretic interiors of the
disjoint open sets $\Etil,K_1,K_2,\dots,U_\infty$, and the perimeters add up
along their shared interfaces. Concretely, write $F_\infty:=U_\infty$,
$F_0:=\Etil$, $F_i:=K_i$. Each pair of these sets is separated by the level
set $\set{u=\hat v}\cap\Om$ or sits across $\partial\Etil$, and by the locality
and structure of the reduced boundary \cite[Thm.~16.3, \S16.2]{Maggi2012} the
essential boundary of $\Etil$ decomposes $\HH^1$-a.e.\ as
\begin{equation}\label{eq:fill-boundary}
   \partial^*\Etil\ =\ \bigl(\partial^*\Etil\cap\partial^*U_\infty\bigr)
   \ \sqcup\ \bigsqcup_i\bigl(\partial^*\Etil\cap\partial^*K_i\bigr)
   \qquad(\HH^1\text{-a.e.}),
\end{equation}
since every point of $\partial^*\Etil$ has, on its outer side, density-one
membership in exactly one of the complementary components $U_\infty,K_1,K_2,
\dots$. Filling the bounded components $K_i$ \emph{removes} from the boundary
exactly the interfaces $\partial^*\Etil\cap\partial^*K_i$ (now interior to
$E^h$) and keeps the outer interface $\partial^*\Etil\cap\partial^*U_\infty=
\partial^*E^h$; moreover, since each $K_i$ has \emph{all} of its essential
boundary shared with $\Etil$ or with other filled components (its complement
in $\R^2$ is $\Etil\cup U_\infty\cup\bigcup_{j\ne i}K_j$, whose density-one
points exhaust $\partial^*K_i$), filling adds no new reduced boundary. Hence
$\partial^*E^h=\partial^*\Etil\cap\partial^*U_\infty=\partial^*\Etil
\setminus\bigsqcup_i\partial^*K_i$ ($\HH^1$-a.e.), and by additivity of
$\HH^1$ over \eqref{eq:fill-boundary},
\begin{equation}\label{eq:fill-perimexact}
   \Per(E^h)=\HH^1(\partial^*E^h)
   =\HH^1(\partial^*\Etil)-\sum_i\HH^1(\partial^*\Etil\cap\partial^*K_i)
   =\Per(\Etil)-\sum_i\Per(K_i),
\end{equation}
the last equality because $\partial^*K_i\subset\partial^*\Etil$ $\HH^1$-a.e.\
(every essential-boundary point of $K_i$ adjacent to $\Etil$ is one of
$\Etil$ adjacent to $K_i$; the components $K_i$ do not touch each other or
$U_\infty$ across a set of positive $\HH^1$-measure, being separated by
$\overline\Etil$). In particular $\Per(E^h)\le\Per(\Etil)$, which is
\eqref{eq:fill-per}.

\emph{\ref{it:fill-key}.} Let the filled holes $K_i$ have areas $A_i^h:=|K_i|$
and perimeters $p_i:=\Per(K_i)$, so $H=\sum_iA_i^h$. By \ref{it:fill-per},
filling removes exactly the boundaries of the holes, $\Per(E^h)=\Per(\Etil)
-\sum_i p_i$ (eq.~\eqref{eq:fill-perimexact}).
By the planar isoperimetric inequality (\cref{thm:isoperimetric}) applied to
each hole, $p_i\ge2\sqrt{\pi A_i^h}$; and by superadditivity of
$t\mapsto\sqrt t$ on non-negative terms,
\begin{equation}\label{eq:fill-superadd}
   \sum_i p_i\ \ge\ 2\sqrt\pi\sum_i\sqrt{A_i^h}\ \ge\ 2\sqrt\pi
   \sqrt{\textstyle\sum_iA_i^h}\ =\ 2\sqrt{\pi H}.
\end{equation}
On the other hand the isoperimetric inequality for $E^h$ itself, of area
$\mu_E+H$, gives $\Per(E^h)\ge2\sqrt{\pi(\mu_E+H)}$. Substituting both into
\eqref{eq:fill-perimexact},
\[
   2\sqrt{\pi(\mu_E+H)}\ \le\ \Per(E^h)\ =\ \Per(\Etil)-\sum_i p_i
   \ \le\ \Per(\Etil)-2\sqrt{\pi H},
\]
so that
\[
   \Per(\Etil)\ \ge\ 2\sqrt\pi\Bigl(\sqrt{\mu_E+H}+\sqrt{H}\Bigr).
\]
Squaring and using $\sqrt{(\mu_E+H)H}\ge\sqrt{\mu_E H}$,
\[
   \Per(\Etil)^2\ \ge\ 4\pi\Bigl(\mu_E+2H+2\sqrt{H(\mu_E+H)}\Bigr)
   \ \ge\ 4\pi\mu_E+8\pi\sqrt{H\mu_E},
\]
whence $D_I(\Etil)=\Per(\Etil)^2-4\pi\mu_E\ge8\pi\sqrt{H\mu_E}$, i.e.
\[
   H\ \le\ \frac{D_I(\Etil)^2}{64\pi^2\mu_E},
\]
which is \eqref{eq:fill-key}; the second inequality there uses
$\mu_E\ge\pi/4$ (\cref{lem:spill}), giving the absolute constant
$C=1/(16\pi^3)$.

\emph{\ref{it:fill-scales}.} By \eqref{eq:diso-inset} and its proof,
$D_I(\Etil)\le D_I(\hat v)+4\pi S\le2D_I(\hat v)\le2A_0\,\dT/\tau$, so
\eqref{eq:fill-key} gives $H\le C(\dT/\tau)^2$, and at $\tau=\sqrt\dT$,
$H\le C\dT$, of the same higher order as the spillover $S$ of
\cref{lem:spill}. The consequences follow:

\noindent\emph{Closeness and radius.} Since $\Etil\subset E^h$ and
$\Etil\subset\Om$ (so $|\Etil\setminus\Om|=0$), we split the symmetric
difference: $|\Om\setminus E^h|\le|\Om\setminus\Etil|$ (as $\Etil\subset E^h$
shrinks the part of $\Om$ outside), while $|E^h\setminus\Om|\le|E^h\setminus
\Etil|+|\Etil\setminus\Om|=H+0=H$. Hence
\[
   |\Om\triangle E^h|=|\Om\setminus E^h|+|E^h\setminus\Om|
   \ \le\ |\Om\setminus\Etil|+H=|\Om\triangle\Etil|+H
   \ \le\ C\sqrt\dT+C\dT\le C\sqrt\dT,
\]
which is \eqref{eq:fill-close}; no assertion that $E^h\subset\Om$ is needed.
The radius obeys $\hat r^h=\sqrt{(\mu_E+H)/\pi}=
\hat r\sqrt{1+H/\mu_E}=\hat r+O(H/\hat r)=\hat r+O(\dT)$ by $H\le C\dT$ and
$\hat r\ge1/2$.

\noindent\emph{Isoperimetric deficit.} Since $\Per(E^h)\le\Per(\Etil)$
\eqref{eq:fill-per} and $|E^h|=\mu_E+H\ge\mu_E$,
\[
   \diso(E^h)=\frac{\Per(E^h)}{2\sqrt{\pi|E^h|}}-1
   \ \le\ \frac{\Per(\Etil)}{2\sqrt{\pi\mu_E}}-1=\diso(\Etil)\le C\,\tfrac\dT\tau,
\]
which is the first bound in \eqref{eq:fill-diso}; and $\Exc(E^h)=
\Per(E^h)-2\pi\hat r^h\le\Per(\Etil)-2\pi\hat r+2\pi(\hat r-\hat r^h)\le
\Exc(\Etil)+CH\le C\,\dT/\tau$, using $|\hat r-\hat r^h|\le CH$.

\noindent\emph{Ball distance.} For any common centre,
$|E^h\triangle B_{\hat r^h}(z)|\le|\Etil\triangle B_{\hat r}(z)|+|E^h
\setminus\Etil|+|B_{\hat r^h}(z)\triangle B_{\hat r}(z)|\le a(\Etil)+H+
C|\hat r^h-\hat r|\le a(\Etil)+CH$; minimising over $z$ gives $a(E^h)\le
a(\Etil)+CH\le a(\Etil)+C\dT$, which is \eqref{eq:fill-a}, and the parenthetical
rate follows from $a(\Etil)\le C(\dT/\tau)^{1/2}$ and $\dT\le(\dT/\tau)^{1/2}$
(as $\dT/\tau\le1$).

\noindent\emph{Scale-invariant deficit.} Torsional rigidity is monotone under
set inclusion: from the variational characterisation
$T(F)=\max\set{2\int_F v-\int_F|\nabla v|^2:\ v\in H^1_0(F)}$ \textup{(}attained
at the torsion function of $F$\textup{)} and the inclusion $H^1_0(\Etil)\subset
H^1_0(E^h)$ (extension by zero, valid as $\Etil\subset E^h$ are open), one has
$T(E^h)\ge T(\Etil)$ for $\Etil\subset E^h$. Hence, with $|E^h|=\mu_E+H$,
\[
   1-\dstar(E^h)=\frac{8\pi T(E^h)}{|E^h|^2}
   \ \ge\ \frac{8\pi T(\Etil)}{(\mu_E+H)^2}
   =\frac{1-\dstar(\Etil)}{(1+H/\mu_E)^2},
\]
so
\[
   \dstar(E^h)\ \le\ 1-\frac{1-\dstar(\Etil)}{(1+H/\mu_E)^2}
   \ \le\ \dstar(\Etil)+\Bigl(1-\frac1{(1+H/\mu_E)^2}\Bigr)
   \ \le\ \dstar(\Etil)+\frac{2H}{\mu_E},
\]
using $1-(1+t)^{-2}\le2t$ for $t\ge0$. With $\mu_E\ge\pi/4$ and $H\le C\dT$
this is $\dstar(E^h)\le\dstar(\Etil)+C\dT\le C\dstar(\Om)$ by
\eqref{eq:dstar-inset} and $\dT=\tfrac\pi8\dstar(\Om)$, which is
\eqref{eq:fill-dstar}.
\end{proof}

\begin{remark}[The hole-filled inset is the object of
\cref{sec:annulus,sec:fold}]\label{rem:fill-use}
\Cref{lem:fill} resolves, unconditionally, the global exterior-connectedness
property required by the trapping annulus and the inscribed-ball lemma: the
hole-filled inset $E^h$ is connected and satisfies $\R^2\setminus
\overline{E^h}$ connected, with all the inset scales of
\cref{prop:insetcontrols} preserved up to an additive $O(\dT)$, of the same
higher order as the spillover. From \cref{sec:annulus} onward the geometric
reduction is therefore applied to $E^h$ in place of $\Etil$; the only change
in any estimate is the replacement of $\mu_E$ by $\mu_E+H=\mu_E+O(\dT)$ and
of $\hat r,\diso,\Exc,a,\dstar$ by their $E^h$-counterparts, with identical
rates. The mass $H$ added by filling, like the spillover $S$ subtracted by
selecting the largest component, is quadratic in the isoperimetric deficit
and hence negligible at the sharp scale.
\end{remark}

\begin{remark}[Self-contained deficit transfer; no forward dependence]
\label{rem:fwd}
The bound \eqref{eq:dstar-inset} is proved above without any appeal to the
cutting principle \cref{lem:cutdeficit} of \cref{sec:cut}: the passage from
$E_{\hat v}$ to its largest component $\Etil$ is handled by the exact
additivity of torsion over components \eqref{eq:Tadd} and the Saint--Venant
bound \eqref{eq:Tspill} on the spillover torsion, both internal to
\cref{sec:identity,sec:inset}. There is therefore \emph{no} circular
dependence on \cref{sec:cut} in this section. (One could alternatively invoke
\cref{lem:cutdeficit}; but its precise hypothesis --- that the excised set lie
in a thin concentric collar of $\Etil$ --- is \emph{not} met by the spillover,
whose components are in general not contained in the trapping annulus of
$\Etil$. The additivity argument avoids this issue entirely, which is why we
prefer it.) For the record, \cref{lem:cutdeficit} itself rests only on
Green-function monotonicity and a boundary-layer barrier, so even the optional
route would be acyclic.
\end{remark}

\begin{remark}[The scales, recorded]\label{rem:scales}
With $\tau=\sqrt\dT$ (\cref{thm:reduction}) the inset $\Etil$ satisfies
$\diso=\diso(\Etil)\le C\sqrt\dT$ and $\Exc=2\pi\hat r\,\diso\le C\sqrt\dT$.
Feeding $\diso$ into the sharp quantitative isoperimetric inequality
(\cref{thm:fmp}) gives the ball-distance
$a:=|\Etil\triangle B_{\hat r}(z)|\le C_{\mathrm F}\sqrt{\diso}\,\mu_E\le
C\,\dT^{1/4}$ at the optimal centre $z$. The trapping width then satisfies
$\eta\le C(\Exc+a)^{1/2}\le C\,\dT^{1/8}$ (\cref{thm:annulus}), an
\emph{absolute} small amplitude. These are the inputs to the annulus and
fold sections; the deficit $\dstar(\Etil)\le C\dstar(\Om)$ and the closeness
$|\Om\triangle\Etil|\le C\sqrt\dT$ are the inputs to the conclusion. We
emphasise that closeness is carried by the \emph{level} $\hat v\sim\sqrt\dT$
through \cref{lem:close}, not by the amplitude $\eta$; and that the deficit
transfer \cref{lem:transfer} delivers $\dstar(\Etil)\lesssim\dstar(\Om)$
without recourse to the inflated single-level deficits.
\end{remark}
